\documentclass[11pt]{article}
\usepackage[margin=1in]{geometry}\usepackage[T1]{fontenc}\usepackage[utf8]{inputenc}
\usepackage{lmodern,amsmath,amssymb,amsthm,mathtools,booktabs,microtype}
\usepackage[colorlinks=true,linkcolor=blue,citecolor=blue,urlcolor=blue]{hyperref}
\hypersetup{pdftitle={A full-support path-variation prior for Bayesian ReLU networks},pdfauthor={ChatGPT-6 Astra; Claude Opus 5.5; Kaarel Hänni; Dmitry Vaintrob}}
\newtheorem{theorem}{Theorem}[section]\newtheorem{lemma}[theorem]{Lemma}
\newtheorem{corollary}[theorem]{Corollary}\newtheorem{proposition}[theorem]{Proposition}
\theoremstyle{definition}\newtheorem{definition}[theorem]{Definition}
\theoremstyle{remark}\newtheorem{remark}[theorem]{Remark}
\newcommand{\R}{\mathbb R}\newcommand{\E}{\mathbb E}\newcommand{\Prob}{\mathbb P}
\newcommand{\Rad}{\mathfrak R}\newcommand{\KL}{\operatorname{KL}}
\newcommand{\op}{\mathrm{op}}\newcommand{\VarPath}{\mathcal V}
\title{A full-support path-variation prior\\for Bayesian ReLU networks}
\author{ChatGPT-6 Astra \and Claude Opus 5.5 \and Kaarel H\"anni \and Dmitry Vaintrob\thanks{Authors are listed in alphabetical order.}}
\date{17 November 2026}
\begin{document}\maketitle
\begin{abstract}
We construct a Bayesian ReLU learner with a full-support prior on one fixed dense architecture and an agnostic oracle controlled by path variation. Start with independent centered Gaussian coordinates and apply a soft penalty for internal amplification and its product with endpoint quadratic energy. No parameters are excluded by a norm constraint. Conditional on the internal matrices, the endpoint coordinates remain Gaussian with gain-normalized variances. The precision-normalized negative log probabilities of function neighborhoods have two-sided bounds in terms of minimal approximate path variation, with no additional learner-width factor. Inputs lie in a Euclidean unit ball and labels may follow any bounded IID law. A likelihood selected from observed free energy gives the ordinary raw posterior mean population loss at most the comparison loss plus $\widetilde O((1+P)\sqrt{HR_*/n}+(1+P)^2H/n)$. Here $P$ is teacher path variation, $H$ hidden depth, and $R_*$ comparison risk. A gain-controlled width-$C$ teacher has $P=O(C)$, so quadratic width and linear depth suffice. The conclusion holds at every learner width above the teacher width and every internal exponent $\gamma\ge1$. A second version keeps the usual layerwise Gaussian base scales and gives a quadratic-width, cubic-depth bound. Its prior is a softly reweighted Gaussian, not the unchanged independent Gaussian tuple. An ordinary-spectral-gain variant retains cancellation between signed weights and gives an oracle with only logarithmic depth cost at fixed input dimension, at the price of dimension-dependent width and accuracy exponents. There is no distribution over network sizes.
\end{abstract}

\section{A function prior without excluding parameter values}
A Bayesian approximation guarantee should compare with a small network even when no small network exactly fits the data. This paper constructs such a learner with a full-support prior on a single dense architecture. Its complexity measure is path variation, an established rescaling-invariant measure for ReLU networks.

Internal matrices may have arbitrarily large norms. A soft amplification penalty reduces prior density when their combined gain amplifies endpoint quadratic energy. Equivalently, the endpoint conditional variances decrease with internal gain. This makes the endpoint energy control path variation while retaining an exact chi-square distribution for that energy. A quantitative Gaussian comparison then proves both a function-prior statement and a sample-complexity oracle.

The prior differs from independent Gaussian coordinates by one explicit soft penalty. We state the normalization explicitly. It is independent of the teacher, target accuracy, and sample, and it imposes no bound on path variation or on coefficient size. Only the likelihood scale is selected from the data.

There are two prior scalings. The main construction matches endpoint precision to the full parameter volume. Section~\ref{sec:original-scale-variant} keeps the usual independent Gaussian base scales and uses a selective comparison. For fixed input dimension and endpoint coefficient bounds, and teacher path variation $P=O(C)$, their sample guarantees are
\[
 \begin{array}{c|c}
 \text{Gaussian base scaling}&\text{sufficient samples, up to logarithms}\\\hline
 \text{endpoint precision matched to parameter volume}
    &C^2H(R_*/\varepsilon^2+1/\varepsilon)\\
 \text{usual layerwise scales, as in~\eqref{eq:usual-base-scales}}
    &C^2H^3(R_*/\varepsilon^2+1/\varepsilon).
 \end{array}
\]
Both priors include an explicitly stated soft gain penalty. Neither bound is claimed for the unchanged independent Gaussian prior.

\section{Architecture, input convention, and path variation}
Let $\mathcal X$ be a compact subset of the Euclidean unit ball in $\R^d$, $d\ge1$. There are $H\ge1$ hidden ReLU layers, a scalar linear readout, and common hidden width $N$:
\[
 h^1(x)=\phi(W^1x),\qquad h^\ell(x)=\phi(W^\ell h^{\ell-1}(x)),\quad
 z_w(x)=v^{\mathsf T}h^H(x),\qquad\phi(t)=\max\{t,0\}.
\]
The first matrix is $N$ by $d$ and the internal matrices are $N$ by $N$. A constant input coordinate may be included in $x$. For ordinary first-layer biases on a unit-ball input $x_0$, use $x=(1,x_0)/\sqrt2$ and count the extra coordinate in $d$. There are no internal biases. Total affine depth is $L=H+1$.

For a network of any finite width, let $a^1_i=\|W^1_{i,:}\|_2$. Its mixed path variation is
\begin{equation}
 P(w)=|v|^{\mathsf T}|W^H|\cdots|W^2|a^1.
 \label{eq:pathnorm}
\end{equation}
For $H=1$, the intervening product is empty. Absolute values on matrices and vectors are entrywise. The Euclidean norm at the first layer matches the input unit ball. In particular $|z_w(x)|\le P(w)$ on $\mathcal X$.

For a function $f$ realizable at depth $H$, define $\VarPath_H(f)$ as the infimum of $P(w)$ over all finite-width depth-$H$ implementations of $f$ on $\mathcal X$. For a width bound $N$ and accuracy $\delta\ge0$, define
\begin{equation}
 \VarPath_{H,N,\delta}(f)=
 \inf\{P(w):\operatorname{width}(w)\le N,\ \|z_w-f\|_{\infty,\mathcal X}\le\delta\}.
 \label{eq:approx-variation}
\end{equation}
The infimum of an empty set is infinity. These are implementation complexity measures, not probabilities.

\section{The full-support Gaussian prior with a soft gain penalty}
For internal matrices define
\[
 g_\ell(W)=\max\{1,\||W^\ell|\|_{\op}\},\qquad
 G(W)=\prod_{\ell=2}^H g_\ell(W).
\]
For $H=1$ the empty product is one. For any $\gamma\ge1$, put
\[
 \omega=N(d+1)+(H-1)N^{\gamma+1}.
\]
Let $\pi^0_{N,\gamma}$ have independent centered Gaussian coordinates, with first-layer variance $N/\omega$, internal variance $N^{-\gamma}$, and readout variance $1/(\omega N)$. Define raw endpoint energies
\[
 u_N=\frac1N\|W^1\|_F^2,\qquad q_N=N\|v\|_2^2.
\]
Our prior is the soft reweighting
\begin{equation}
 \pi_{N,\gamma}(dw)=\frac1{Z_0}
 \exp\left[-\frac\omega2\bigl\{(G(W)-1)(u_N+q_N)+(G(W)^2-1)\bigr\}\right]\pi^0_{N,\gamma}(dw).
 \label{eq:prior}
\end{equation}
The normalizer satisfies $0<Z_0\le1$. This density is positive at every finite parameter vector. It mixes over neither widths nor depths.

Define the gain-normalized endpoint energy and full normalized energy by
\begin{align}
 Q_N&=G(u_N+q_N),\notag\\
 E_N(w)&=Q_N(w)+G(W)^2-1+\frac{N^\gamma}{\omega}\sum_{\ell=2}^H\|W^\ell\|_F^2.
 \label{eq:energy}
\end{align}
The parameter density is proportional to $e^{-\omega E_N(w)/2}$. The terms $(G-1)(u_N+q_N)$ and $G^2-1$ penalize amplification rather than enforcing an internal norm bound. The second term gives a width-independent moment bound for internal gain; the first links endpoint energy to function complexity.

Conditional on the internal matrices, the first layer and readout are independent centered Gaussians with variances
\begin{equation}
 \operatorname{Var}(W^1_{ji}\mid W^2,\ldots,W^H)=\frac{N}{\omega G},\qquad
 \operatorname{Var}(v_j\mid W^2,\ldots,W^H)=\frac1{\omega NG}.
 \label{eq:conditional-prior}
\end{equation}
The internal marginal is its independent Gaussian base tilted by $G^{-N(d+1)/2}e^{-\omega(G^2-1)/2}$, divided by $Z_0$. These conditional identities account for every normalization factor in the proof.

A comparison with width $C$ needs only $N\ge C$, without an input-dimension threshold or a limiting operation.

\section{Path-flow balancing and the exact energy characterization}
\begin{lemma}[Path-flow balancing]
\label{lem:flow}
Every finite ReLU network with path variation $P>0$ has an equivalent implementation at the same width and depth such that every internal absolute spectral norm is at most one and
\begin{equation}
 \|W^1\|_F^2=P,\qquad \|v\|_2^2=P.
 \label{eq:balanced-flow}
\end{equation}
The path variation is unchanged. If $P=0$, the function is zero.
\end{lemma}
\begin{proof}
Write $A^\ell=|W^\ell|$. Starting from $a^1_i=\|W^1_{i,:}\|_2$, set $a^\ell=A^\ell a^{\ell-1}$. Starting from $b^H=|v|$, set $b^{\ell-1}=(A^\ell)^{\mathsf T}b^\ell$. At every hidden layer,
\[
 \sum_j a^\ell_jb^\ell_j=P.
\]
Neurons with zero forward or backward mass lie on no nonzero input-to-output path and can be disconnected without changing the function. On the remaining neurons define
\[
 s^\ell_j=\sqrt{b^\ell_j/a^\ell_j},\qquad
 r^\ell_j=\sqrt{a^\ell_jb^\ell_j}.
\]
Multiply hidden activation $j$ in layer $\ell$ by $s^\ell_j$ by rescaling its incoming weights and inversely rescaling its outgoing weights. ReLU positive homogeneity preserves the exact function. The transformed internal absolute matrix is
$\widetilde A^\ell=\operatorname{diag}(s^\ell)A^\ell\operatorname{diag}(s^{\ell-1})^{-1}$, and it satisfies
\[
 \widetilde A^\ell r^{\ell-1}=r^\ell,\qquad
 (\widetilde A^\ell)^{\mathsf T}r^\ell=r^{\ell-1}.
\]
Weighted Cauchy--Schwarz gives, for any vector $x$ on the active coordinates,
\begin{align*}
 \sum_j\left(\sum_i\widetilde A^\ell_{ji}x_i\right)^2
 &\le\sum_j r^\ell_j\sum_i\widetilde A^\ell_{ji}\frac{x_i^2}{r^{\ell-1}_i}\\
 &=\sum_i\frac{x_i^2}{r^{\ell-1}_i}
                ((\widetilde A^\ell)^{\mathsf T}r^\ell)_i
 =\sum_i x_i^2.
\end{align*}
Thus its operator norm is at most one. The transformed first-row Euclidean norm is $s^1_ja^1_j=r^1_j$, and the transformed readout magnitude is $b^H_j/s^H_j=r^H_j$. Their squared sums are both $P$, proving~\eqref{eq:balanced-flow}. Positive neuron rescaling cancels on every path, so path variation is preserved. If $P=0$, every absolute path product is zero, and $|z_w|\le P$ gives the zero function.
\end{proof}

\begin{lemma}[Cloning without a width divisibility condition]
\label{lem:clone}
A teacher of width $C$ and path variation $P$ has an exact implementation at every $N\ge C$ with
\begin{equation}
 Q_N=2P,\qquad
 \sum_{\ell=2}^H\|W^\ell\|_F^2\le(H-1)C.
 \label{eq:clone}
\end{equation}
For the zero function, the all-zero implementation suffices.
\end{lemma}
\begin{proof}
First apply Lemma~\ref{lem:flow}. Put $k=\lfloor N/C\rfloor$, clone every teacher neuron exactly $k$ times, and leave the other neurons disconnected with zero weights. The internal blocks are $W^\ell\otimes(J_k/k)$. Both their absolute spectral norms and their squared Frobenius norms are unchanged. Since a $C$ by $C$ matrix of operator norm at most one has squared Frobenius norm at most $C$, the internal bound follows.

The endpoint energies before the final balancing are $u_N=(k/N)P$ and $q_N=(N/k)P$. Multiply the first layer by $\sqrt{N/k}$ and divide the readout by the same factor. Positive homogeneity again preserves the function, and now $u_N=q_N=P$. This proves~\eqref{eq:clone} exactly, including widths not divisible by $C$.
\end{proof}

\begin{theorem}[An exact variational path-prior characterization]
\label{thm:energy-characterization}
For every network,
\begin{equation}
 E_N(w)\ge Q_N(w)\ge2P(w).
 \label{eq:energy-lower}
\end{equation}
Conversely, a teacher of width $C$ and path variation $P$ has, at each $N\ge C$, an exact implementation with
\begin{equation}
 E_N\le2P+\frac{N^\gamma(H-1)C}{\omega}\le2P+\frac CN.
 \label{eq:energy-upper}
\end{equation}
Consequently, for every finite-network-realizable function $f$ at depth $H$,
\begin{equation}
 \lim_{N\to\infty}\inf_{\substack{w\text{ of width }N\\z_w=f\text{ on }\mathcal X}}E_N(w)
 =2\VarPath_H(f).
 \label{eq:energy-limit}
\end{equation}
\end{theorem}
\begin{proof}
By the definition of $G$,
\[
 P(w)\le G\|v\|_2\left(\sum_i\|W^1_{i,:}\|_2^2\right)^{1/2}
          =G\sqrt{u_Nq_N}\le Q_N/2.
\]
This proves~\eqref{eq:energy-lower}; Lemma~\ref{lem:clone} proves~\eqref{eq:energy-upper}. The lower bound in~\eqref{eq:energy-limit} holds at every width. For any $\epsilon>0$, choose a finite-width implementation of $f$ with path variation at most $\VarPath_H(f)+\epsilon$. Its width is fixed while $N$ grows, so~\eqref{eq:energy-upper} gives upper limit at most $2\VarPath_H(f)+2\epsilon$. Let $\epsilon$ decrease to zero.
\end{proof}

\section{A full-support Gaussian comparison}
Let $D_{\mathrm{int}}=(H-1)N^2$ and $D_{\mathrm{end}}=N(d+1)$. The internal marginal of the prior has density
\begin{equation}
 \nu(dX)\ \propto\
 \exp\left[-\frac12\sum_{\ell=2}^HN^\gamma\|X^\ell\|_F^2\right]
 G(X)^{-D_{\mathrm{end}}/2}
 \exp\left[-\frac\omega2\{G(X)^2-1\}\right]dX.
 \label{eq:internal-marginal}
\end{equation}
For $H=1$ there are no internal coordinates and $G=1$.

\begin{lemma}[A width-independent internal gain moment]
\label{lem:gain-moment}
For $0<t<1$, the internal prior satisfies
\begin{equation}
 \E_\nu\exp\left[\frac{\omega t}{2}(G^2-1)\right]
 \le(1-t)^{-D_{\mathrm{int}}/2}.
 \label{eq:gain-mgf}
\end{equation}
Consequently $\E_\nu G^2\le1+D_{\mathrm{int}}/\omega\le2$.
\end{lemma}
\begin{proof}
For $0<c\le1$, each factor in $G(cX)$ is $\max\{1,c\||X^\ell|\|_{\op}\}$. If none exceeds one, $G(cX)=1$. Otherwise their product is at most $cG(X)$: if $k\ge1$ factors exceed one, its factor $c^k$ is at most $c$, and every omitted factor in $G(X)$ is at least one. Therefore
\[
 G(cX)\le\max\{1,cG(X)\},\qquad
 G(cX)^2-1\le c^2(G(X)^2-1).
\]
The Gaussian quadratic potential and $\log G$ also decrease under this scaling. Put $c=\sqrt{1-t}$ and change variables $Y=cX$ in the normalizing integral of~\eqref{eq:internal-marginal}. Its Jacobian is $c^{D_{\mathrm{int}}}$, and the preceding inequalities lower-bound the transformed integral by the original normalizer times
$c^{D_{\mathrm{int}}}\E_\nu e^{\omega t(G^2-1)/2}$. This proves~\eqref{eq:gain-mgf}.

Jensen gives
$\omega t\E(G^2-1)/2\le-(D_{\mathrm{int}}/2)\log(1-t)$.
Divide by $t$ and let $t$ decrease to zero. The moment is finite by~\eqref{eq:gain-mgf}, and $D_{\mathrm{int}}\le\omega$, proving the claim.
\end{proof}

\begin{lemma}[A uniform function tube and its entropy]
\label{lem:comparison}
Let $g$ have width $C$, depth $H$, and mixed path variation $P$. At every $N\ge C$, every $\gamma\ge1$, and every $0<\eta<1$, there is a law $\rho_{g,\eta}$ relative to the full-support prior such that
\begin{align}
 \frac2\omega\KL(\rho_{g,\eta}\|\pi_{N,\gamma})
 &\le3P+1+2\log(1/\eta),\label{eq:comparison-entropy}\\
 \left(\int\|z_w-g\|_{\infty,\mathcal X}^2d\rho_{g,\eta}\right)^{1/2}
 &\le6(H+1)(P+1)\eta.\label{eq:comparison-uniform}
\end{align}
\end{lemma}
\begin{proof}
Use the exact implementation from Lemma~\ref{lem:clone}, with raw endpoint energies $u_N=q_N=P$ and internal matrices $A^\ell$ satisfying $\||A^\ell|\|_{\op}\le1$. Draw the whole internal tuple $X$ from its actual marginal~\eqref{eq:internal-marginal}, and set
\[
 W^\ell=T_\eta X^\ell=(1-\eta)A^\ell+\eta X^\ell.
\]
Conditional on this transformed tuple, draw the first layer and readout independently with the cloned means and $\eta^2$ times the respective conditional prior covariances.

Convexity of $W\mapsto\max\{1,\||W|\|_{\op}\}$ gives
\[
 g_\ell(T_\eta X)\le1-\eta+\eta g_\ell(X)\le g_\ell(X).
\]
Thus $G(T_\eta X)\le G(X)$. In fact,
$G(T_\eta X)\le1+\eta(G(X)-1)$, by expanding the product of $1+\eta(g_\ell(X)-1)$ and using $\eta^j\le\eta$ for $j\ge1$.

The density~\eqref{eq:internal-marginal} is centrally symmetric, so $\E X=0$. Its affine Jacobian is $\eta^{D_{\mathrm{int}}}$. Writing the exact log-density ratio, the internal relative entropy is
\begin{align*}
 D_{\mathrm{int}}\log(1/\eta)
 &+\frac12\sum_{\ell=2}^HN^\gamma\left[(1-\eta)^2\|A^\ell\|_F^2
                      -(1-\eta^2)\E\|X^\ell\|_F^2\right]\\
 &+\frac{D_{\mathrm{end}}}{2}\E\log\frac{G(T_\eta X)}{G(X)}
  +\frac\omega2\E\left[G(T_\eta X)^2-G(X)^2\right].
\end{align*}
The last two terms are nonpositive. After normalization by $2/\omega$, the mean cost is at most
$N^\gamma(H-1)C/\omega\le C/N\le1$.

The conditional endpoint Gaussian entropy has normalized mean cost
$2P\E G(T_\eta X)\le2\sqrt2P\le3P$, by Lemma~\ref{lem:gain-moment}. Its covariance cost is at most $2D_{\mathrm{end}}\log(1/\eta)/\omega$. Since
\[
 \frac{D_{\mathrm{int}}+D_{\mathrm{end}}}{\omega}
 =\frac{(H-1)N^2+N(d+1)}{(H-1)N^{\gamma+1}+N(d+1)}\le1,
\]
the chain rule proves~\eqref{eq:comparison-entropy}. All normalizers cancel in the exact marginal or conditional density ratios.

For the uniform error, first hold the internal tuple fixed. Write $G_T=G(T_\eta X)$ and $M_\ell=\|T_\eta X^\ell\|_{\op}\le g_\ell(X)$. The propagated first-layer noise, with hidden norms divided by $\sqrt N$, has squared expectation at most
\[
 \E\left[\frac{\prod_{\ell=2}^HM_\ell^2}{N}
                  \|W^1-W^{1,*}\|_F^2\right]
 \le\eta^2\frac{Nd}{\omega}\E G_T\le\sqrt2\eta^2.
\]
Here the conditional first-layer covariance contains $G_T^{-1}$ and $\prod M_\ell\le G_T$. This estimate handles its dependence on the whole internal tuple.

For internal perturbations,
$\|T_\eta X^\ell-A^\ell\|_{\op}\le2\eta g_\ell(X)$.
The teacher hidden states have norm at most $\sqrt{NP}$. Telescoping the compositions, each layer's normalized error is at most
$2\eta\sqrt P\,g_\ell(X)\prod_{j>\ell}g_j(X)\le2\eta\sqrt P\,G(X)$.
Summing and using $\E G^2\le2$ gives root mean squared normalized error at most $3(H-1)\eta\sqrt P$. No independence between internal layers is required. The cloned readout therefore contributes at most
$2\eta\sqrt P+3(H-1)\eta P$ to output error.

Conditional on internals, the readout noise is independent of the first layer, has expected squared norm $\eta^2/(\omega G_T)$, and every hidden-state norm is bounded by $G_T\|W^1\|_F$. Thus
\begin{align*}
 \E\sup_x|(v-v^*)^{\mathsf T}h^H(x)|^2
 &\le\eta^2\frac{NP}{\omega}\E G_T
             +\eta^4\frac{N^2d}{\omega^2}\\
 &\le\eta^2(2P+1).
\end{align*}
Minkowski and elementary square-root bounds now give~\eqref{eq:comparison-uniform} with the stated slack.
\end{proof}

For learning, set
\begin{equation}
 \eta_n=\frac1{64n(H+1)(P+1)^2},\qquad
 T_g=6+3P+2\log(1/\eta_n).
 \label{eq:teacher-budget}
\end{equation}
The uniform error in~\eqref{eq:comparison-uniform} is at most $3/[32n(P+1)]$. Because $|g|\le P$ and $|Y|\le1$, the expected empirical squared-loss increase is at most $1/(2n)$. The Gibbs variational formula consequently gives
\begin{equation}
 F_S(\lambda)\le\widehat R_S(g)+\frac1{2n}+\lambda T_g
 \qquad(\lambda>0).
 \label{eq:free-energy-comparison}
\end{equation}

\section{The induced prior on function neighborhoods}
For a bounded function $f$ on $\mathcal X$ and $0<\delta<1$, write
\[
 B_\delta(f)=\{w:\|z_w-f\|_{\infty,\mathcal X}\le\delta\}.
\]
These events are measurable because the network map is continuous in parameters in the uniform norm on the compact input domain.

\begin{theorem}[Two-sided finite function-prior bounds]
\label{thm:prior}
For every $N\ge1$ and $\gamma\ge1$,
\begin{equation}
 -\frac1\omega\log\pi_{N,\gamma}(B_\delta(f))
 \ge\frac12\VarPath_{H,N,\delta}(f)-\frac{\log2}{2}.
 \label{eq:prior-lower}
\end{equation}
If $p=\VarPath_{H,N,\delta/2}(f)<\infty$, then
\begin{equation}
 -\frac1\omega\log\pi_{N,\gamma}(B_\delta(f))
 \le3p+1+2\log\frac{24(H+1)(p+1)}\delta+2\log2.
 \label{eq:prior-upper}
\end{equation}
In particular the precision-normalized negative log prior is controlled by minimal approximate path variation, with only the displayed logarithmic accuracy and depth terms. No learner-width or input-dimension factor is suppressed.
\end{theorem}
\begin{proof}
Conditional on the internal matrices, the scaled endpoint coordinates are independent standard Gaussians. Their squared sum is $\omega Q_N$, so unconditionally
\begin{equation}
 \omega Q_N\sim\chi^2_{N(d+1)},\qquad
 \frac{N(d+1)}\omega\le1.
 \label{eq:chi-square}
\end{equation}
Any network in $B_\delta(f)$ has $P(w)\ge\VarPath_{H,N,\delta}(f)$ and $Q_N\ge2P(w)$ by~\eqref{eq:energy-lower}. Markov's inequality with
$\E e^{\omega Q_N/4}=2^{N(d+1)/2}$ yields
\[
 \pi(B_\delta(f))\le
 \exp\left[-\frac\omega2\VarPath_{H,N,\delta}(f)
                      +\frac{N(d+1)\log2}{2}\right],
\]
which proves~\eqref{eq:prior-lower}. An infinite infimum means the event is empty.

For the upper bound, choose any width-at-most-$N$ network $g$ within $\delta/2$ of $f$, with path variation $P$. Apply Lemma~\ref{lem:comparison} at
$\eta=\delta/[24(H+1)(P+1)]$. Its uniform-error $L^2$ norm is at most $\delta/4$, so its comparison law gives probability at least $3/4$ to $B_{\delta/2}(g)\subset B_\delta(f)$. The binary relative-entropy inequality implies, for any event $A$ with $\rho(A)\ge1/2$,
\[
 \log\frac1{\pi(A)}\le2\KL(\rho\|\pi)+2\log2.
\]
Use~\eqref{eq:comparison-entropy}, divide by $\omega\ge1$, and let $P$ decrease to the infimum $p$. This proves~\eqref{eq:prior-upper}.
\end{proof}

The different radii in the two bounds are intentional: approximation at half the requested accuracy leaves room for an open neighborhood of the comparison implementation. The theorem is about function neighborhoods, not positive mass on one exact continuous-domain function. It gives bounds on prior preference; it does not assert an exact monotone ordering between every pair of functions.

\section{A dimension-free capacity bound}
Let $\mathcal F_t$ be the union over all positive widths $N$ of networks with $Q_N\le t$, at fixed hidden depth. Put
\[
 K_H=1+\sqrt{2H\log2}.
\]

\begin{lemma}[Capacity and amplitude]
\label{lem:capacity}
On every input sample in $\mathcal X$,
\[
 \widehat\Rad_n(\mathcal F_t)\le\frac{tK_H}{\sqrt n},\qquad
 |z_w(x)|\le P(w)\le Q_N(w)/2.
\]
The same capacity bound holds for worst-case samples, with no dimension dependence.
\end{lemma}
\begin{proof}
The amplitude statement follows from~\eqref{eq:energy-lower}. To bound capacity, normalize every hidden feature by its forward path mass: divide $h_j^\ell$ by $a_j^\ell$, discarding zero-mass coordinates. ReLU homogeneity makes every first-row Euclidean norm at most one, every internal absolute row sum at most one, and the transformed readout's absolute coefficient sum equal to $P(w)\le t/2$.

Let $\mathcal A_0$ be the sample vectors of all Euclidean-unit-norm linear functions. Recursively let $\mathcal A_\ell$ be the closed absolute convex hull of the coordinatewise ReLU images of $\mathcal A_{\ell-1}$. Every normalized hidden feature belongs to its corresponding class. For $M_\ell=\sup_{f\in\mathcal A_\ell}\sum_i\epsilon_if_i$ and $s>0$, exponential scalar contraction gives
\[
 \E e^{sM_\ell}\le2\E e^{sM_{\ell-1}}.
\]
Indeed, absolute convexification gives a maximum of positive and negative suprema. Its exponential is at most their summed exponentials, which have equal expectation. Contraction for either supremum has constant one: conditional on other signs, Lipschitzness controls the sum of the two coordinate suprema, and $|\phi(t)|\le|t|$ controls their maximum; these imply the inequality for the sum of exponentials.

At the initial layer, $M_0=\|\sum_i\epsilon_iX_i\|_2$. Jensen gives $\E M_0\le\sqrt n$. Changing sign $i$ changes $M_0$ by at most $2\|X_i\|_2\le2$, so the bounded-difference exponential inequality gives
\[
 \E e^{sM_0}\le e^{s\sqrt n+s^2n/2}.
\]
This last inequality also follows directly by applying Hoeffding's lemma to the Doob martingale of $M_0$. Iterating the contraction, taking logarithms, and optimizing $s$ yields $\E M_H\le\sqrt nK_H$. Multiply by the readout absolute sum at most $t/2$ and divide by $n$; the displayed capacity bound leaves a harmless factor two of slack. Separability follows from countable dense parameter subclasses at each width and a countable union over widths.
\end{proof}

\section{A data-dependent likelihood with a raw-risk certificate}
Let $R(f)=\E(f(X)-Y)^2$ and $\widehat R_S(f)=n^{-1}\sum_i(f(X_i)-Y_i)^2$. With Gaussian observation variance $n\lambda/\omega$, define
\[
 \Pi_{\lambda,S}(dw)=Z^{-1}e^{-\omega\widehat R_S(z_w)/(2\lambda)}\pi_{N,\gamma}(dw),\qquad
 F_S(\lambda)=-\frac{2\lambda}{\omega}\log Z.
\]
The variational identity is $F_S(\lambda)=\inf_\rho\{\int\widehat R_Sd\rho+2\lambda\KL(\rho\|\pi)/\omega\}$. The centered readout, independent of the first layer conditional on internal matrices, cancels cross terms and
\[
 \E_\pi[z_w(x)^2\mid W^2,\ldots,W^H]\le\frac{G^2}{\omega NG}\frac{N^2}{\omega G}
                 \le\frac N{\omega^2}\le1.
\]
Therefore $0\le F_S(\lambda)\le2$ for every sample and every $\lambda>0$.
We state all constants used in the selection rule. Let $C_*\ge1$ be a universal constant for the smooth-loss theorem of Srebro, Sridharan, and Tewari~\cite[Theorem~1]{sst}. In the form used here, for a nonnegative $s_0$-smooth loss bounded by $b$ on $\mathcal F$, with probability at least $1-\eta$, simultaneously for $f\in\mathcal F$,
\begin{align}
 R(f)-\widehat R_S(f)\le C_*\bigg[&\sqrt{\widehat R_S(f)}
 \left(\sqrt{s_0}\log^{3/2}(2n)\Rad_n(\mathcal F)+\sqrt{\frac{b\log(1/\eta)}n}\right)\notag\\
 &+s_0\log^3(2n)\Rad_n(\mathcal F)^2+\frac{b\log(1/\eta)}n\bigg].
 \label{eq:sst}
\end{align}
The Rademacher complexity here is worst-case empirical complexity. Squared loss is $2$-smooth.

For $0<\delta<1$, put
\begin{align}
 M_2&=1+(2D_0+8)^2, & K_n&=2D_0\log2+4\log(6M_2n^4)+1,\notag\\
 c_0&=\frac18+\frac{D_0\log2}{4}, &
 J_n&=\left\lceil\log_2(4n+K_n+2)\right\rceil,\notag\\
 t_n&=\max\left\{1,\log\frac{4(J_n+1)}\delta\right\},\notag\\
 a_n&=\max\left\{\frac1{\sqrt n},\ 2C_*\left[
       \frac{\sqrt2K_H\log^{3/2}(2n)}{\sqrt n}+2\sqrt{\frac{t_n}n}\right]\right\}.
 \label{eq:stat-constants}
\end{align}
Set $D_0=1$ and recall $K_H$ from~Lemma~\ref{lem:capacity}. Select the unique positive solution of
\begin{equation}
 \widehat\lambda_S^2
   =256a_n^2\{2F_S(\widehat\lambda_S)+\widehat\lambda_S(K_n+1)\}.
 \label{eq:calibration}
\end{equation}
This one-dimensional rule uses the observed free energy, not teacher resources or the population risk. The learner predicts with the posterior at this single selected value.

\begin{lemma}[Existence, uniqueness, and measurability]
\label{lem:root}
The root in~\eqref{eq:calibration} exists, is unique, and is a measurable function of the sample. It satisfies
$\widehat\lambda_S\ge256a_n^2(K_n+1)\ge1/n$.
No upper bound on $\widehat\lambda_S$ is required.
\end{lemma}
\begin{proof}
The likelihood integral makes $F_S$ continuous, and $F_S(\lambda)/\lambda$ is decreasing. The map
\[
 \Psi_S(\lambda)=\frac{\lambda^2}{2F_S(\lambda)+\lambda(K_n+1)}
 =\frac{\lambda}{2F_S(\lambda)/\lambda+K_n+1}
\]
is therefore continuous and strictly increasing. It tends to zero at zero because it is at most $\lambda/(K_n+1)$, and tends to infinity because $F_S\le2$. It crosses $256a_n^2$ exactly once. The lower bound follows directly from the root equation. For each $t>0$, the event $\{\widehat\lambda_S<t\}$ equals $\{\Psi_S(t)>256a_n^2\}$, proving measurability.
\end{proof}

\begin{lemma}[One statistical event for all widths and scales]
\label{lem:stat-event}
With probability at least $1-\delta/2$, simultaneously for every network with $Q_N\le4n+K_n+1$,
\begin{equation}
 R(z_w)-\widehat R_S(z_w)
 \le a_n\sqrt{\widehat R_S(z_w)}(1+Q_N)+a_n^2(1+Q_N)^2.
 \label{eq:stat-event}
\end{equation}
\end{lemma}
\begin{proof}
Apply~\eqref{eq:sst} to $\mathcal F_{2^j}$, $0\le j\le J_n$, with failure probability $\delta/[2(J_n+1)]$. By~\eqref{eq:energy-lower}, its loss is bounded by $4(1+2^j)^2$, and by Lemma~\ref{lem:capacity}, its Rademacher complexity is at most $K_H(1+2^j)/\sqrt n$. Use the unit cutoff for $Q_N\le1$ and the least containing dyadic cutoff otherwise; then $1+2^j\le2(1+Q_N)$. The definition of $a_n$ and a union bound prove the claim. The energy-ball classes already include every width, so the event does not depend on $N$, $\gamma$, or $\lambda$.
\end{proof}

\begin{lemma}[Weighted empirical-objective tails]
\label{lem:tail}
For every sample and $\lambda>0$, the posterior raw-risk integral on
\[
 \{\widehat R_S+\lambda Q_N>2F_S(\lambda)+\lambda K_n\}
\]
is at most $n^{-4}$. Also, for every finite-entropy law $\rho$ relative to the prior,
\begin{equation}
 \int Q_Nd\rho\le\frac4\omega\KL(\rho\|\pi)+2D_0\log2.
 \label{eq:entropy}
\end{equation}
\end{lemma}
\begin{proof}
On $J:=\widehat R_S+\lambda Q_N>u$,
\[
 e^{-\omega\widehat R_S/(2\lambda)}
 \le e^{-\omega u/(4\lambda)}e^{\omega Q_N/4}.
\]
By~\eqref{eq:chi-square}, exponential tilting by $e^{\omega Q_N/4}$ doubles the chi-square scale and contributes the factor $2^{N(d+1)/2}$. Under that tilted law,
\[
 \E Q_N^2=\frac4{\omega^2}N(d+1)\{N(d+1)+2\}\le(2D_0+8)^2.
\]
After dividing by the posterior normalizer,
\[
 \int(1+Q_N^2)\mathbf1_{\{J>u\}}d\Pi_\lambda
 \le M_2\exp\left[\frac{\omega F_S}{2\lambda}-\frac{\omega u}{4\lambda}
                       +\frac{N(d+1)\log2}{2}\right].
\]
Raw risk is at most $2Q_N^2+2\le6(1+Q_N^2)$. Use $u=2F_S+\lambda K_n$, $N(d+1)\le\omega D_0$, and $\omega\ge1$ to obtain the claimed tail bound.

Finally $\log\int e^{\omega Q_N/4}d\pi=N(d+1)\log2/2\le\omega D_0\log2/2$. The entropy variational inequality gives~\eqref{eq:entropy}.
\end{proof}

\begin{theorem}[An observable certificate and an empirical oracle]
\label{thm:certificate}
On the event of Lemma~\ref{lem:stat-event}, simultaneously for every $N\ge1$ and $\gamma\ge1$,
\begin{equation}
 \int R(z_w)d\Pi_{\widehat\lambda_S}
 \le F_S(\widehat\lambda_S)+c_0\widehat\lambda_S+n^{-4}.
 \label{eq:certificate}
\end{equation}
If $F_S(\lambda)\le u+\lambda T$ for every positive $\lambda$, put
$\mathcal B=2T+K_n+1+c_0$. Then
\begin{equation}
 \int R(z_w)d\Pi_{\widehat\lambda_S}
 \le u+16\sqrt2\mathcal B a_n\sqrt u+256\mathcal B^2a_n^2+n^{-4}.
 \label{eq:empirical-oracle}
\end{equation}
Both bounds also hold for the ordinary raw posterior mean.
\end{theorem}
\begin{proof}
Write $\lambda=\widehat\lambda_S$ and $G=2F_S(\lambda)+\lambda(K_n+1)$. Off the tail in Lemma~\ref{lem:tail},
\[
 \widehat R_S\le G,\qquad1+Q_N\le G/\lambda,\qquad Q_N\le4n+K_n.
\]
The last inequality uses $F_S\le2$ and Lemma~\ref{lem:root}. The two terms of~\eqref{eq:stat-event} are bounded, respectively, by
\[
 a_n\sqrt G(1+Q_N)=\frac\lambda{16}(1+Q_N),
 \qquad
 a_n^2\frac G\lambda(1+Q_N)=\frac\lambda{256}(1+Q_N).
\]
Their sum is at most $\lambda(1+Q_N)/8$. Integrate and use~\eqref{eq:entropy}. The energy term costs at most one quarter of the relative-entropy penalty retained in $F_S$, plus $\lambda D_0\log2/4$. Adding $\lambda/8$ and the weighted tail gives~\eqref{eq:certificate}.

For the comparison, the root implies
\[
 \lambda^2\le512a_n^2u+256a_n^2(2T+K_n+1)\lambda.
\]
The elementary implication $x^2\le bx+c\Rightarrow x\le b+\sqrt c$ gives
$\lambda\le256a_n^2\mathcal B+16\sqrt2a_n\sqrt u$.
Substitute into~\eqref{eq:certificate} and use $T+c_0\le\mathcal B$ to obtain~\eqref{eq:empirical-oracle}. All raw moments exist by~\eqref{eq:energy-lower} and the Gaussian endpoint tails. Jensen gives the mean-prediction claim.
\end{proof}

\begin{theorem}[A dimension-free path-variation oracle]
\label{thm:main}
Fix any ReLU teacher $g$ with hidden depth $H$, width $C$, and mixed path variation $P$. Define $T_g$ by~\eqref{eq:teacher-budget}, and put $\mathcal B_g=2T_g+K_n+1+c_0$. With probability at least $1-\delta$, simultaneously for every $N\ge C$ and every $\gamma\ge1$,
\begin{align}
 R\left(\int z_wd\Pi_{\widehat\lambda_S}\right)
 \le\int R(z_w)d\Pi_{\widehat\lambda_S}
 \le R(g)&+25a_n\mathcal B_g\sqrt{R(g)}\notag\\
 &+280a_n^2\mathcal B_g^2+n^{-4}.
 \label{eq:population-oracle}
\end{align}
The calibration uses none of $C,P,R(g)$ and no coefficient bound is required. The data may have any IID law on $\mathcal X\times[-1,1]$.
\end{theorem}
\begin{proof}
The comparison~\eqref{eq:free-energy-comparison} supplies $u=\widehat R_S(g)+1/(2n)$. Put $E=6+3P$. Then $|g|\le P\le E$ and $T_g\ge E\ge4$. Bernstein gives
\[
 \widehat R_S(g)\le R(g)+(E+1)\sqrt{\frac{2R(g)\log(2/\delta)}n}
                  +\frac{2(E+1)^2\log(2/\delta)}{3n}
\]
except on an event of probability at most $\delta/2$. The inequalities $(E+1)^2\le2E^2$, $a_n^2\ge16\log(2/\delta)/n$, and $a_n^2\ge1/n$ imply
$u\le(\sqrt{R(g)}+a_nE)^2$.
Substitute this and its square root into~\eqref{eq:empirical-oracle}, using $E\le T_g\le\mathcal B_g$. The square-root coefficient is at most $2+16\sqrt2<25$, and the quadratic coefficient at most $257+16\sqrt2<280$. Combining the two sample events proves the result. Both are common to every eligible learner width and exponent.
\end{proof}

\begin{corollary}[Sample complexity in path variation]
\label{cor:polynomial}
For a comparison network with path variation $P$ and risk $R_*$, a sufficient sample count for oracle excess at most $\varepsilon\in(0,1)$ is
\begin{equation}
 n=\widetilde O\left(H(1+P)^2\left\{\frac{R_*}{\varepsilon^2}+\frac1\varepsilon\right\}\right).
 \label{eq:sample-budget}
\end{equation}
Only logarithms in $n,H,P,1/\varepsilon,1/\delta$ are suppressed. There is no input-dimension factor once $N\ge C$. The teacher resources and risk are used to state the guarantee, not to choose the posterior.
\end{corollary}
\begin{proof}
Here $a_n^2=\widetilde O(H/n)$ and $\mathcal B_g=\widetilde O(1+P)$. Substitute in~\eqref{eq:population-oracle}. These estimates follow directly from the displayed constants with $D_0=1$ and $K_H=1+\sqrt{2H\log2}$.
\end{proof}

\begin{corollary}[Quadratic width and linear depth under controlled gain]
\label{cor:width-depth}
Suppose a width-$C$ teacher has first-row Euclidean norms at most $B_1$, readout coefficients at most $B_2$, and
\[
 G=\prod_{\ell=2}^H\bigl\||W^\ell|\bigr\|_{\op}.
\]
Then $P\le GCB_1B_2$. In particular, at fixed $B_1,B_2$ and $G\le1$, the sufficient sample count is
\[
 \widetilde O\left(C^2H\left\{\frac{R_*}{\varepsilon^2}+\frac1\varepsilon\right\}\right).
\]
Only the product of gains is needed for this conclusion; the teacher need not have every gain at most one.
\end{corollary}
\begin{proof}
Apply Cauchy--Schwarz to~\eqref{eq:pathnorm}: $P\le\|v\|_2G\|a^1\|_2\le(\sqrt C B_2)G(\sqrt C B_1)$. Then use Corollary~\ref{cor:polynomial}.
\end{proof}

For arbitrary coefficient-bounded teachers, this estimate still gives a fixed-depth polynomial: if every internal coefficient and readout coefficient is at most $B$ in absolute value and first-row norms are at most $B_1$, then $P\le C^HB^HB_1$. Under a coordinatewise first-layer cap $B$ as well, $B_1\le\sqrt d B$, giving $P\le\sqrt d C^HB^{H+1}$. Its possible exponential dependence on depth is not removed by neuron rescaling.

\section{Keeping the usual independent Gaussian base scales}
\label{sec:original-scale-variant}
The preceding prior strengthens the endpoint Gaussian precision to match the full parameter volume. A more selective comparison retains the usual layerwise Gaussian base scales, at the cost of an additional polynomial dependence on depth. It still uses the same soft gain mechanism and excludes no parameters.

Keep the input convention $\mathcal X\subset B_2^d$. Fix $\gamma\ge1$ and put
\[
 \omega=N^\gamma,\qquad D_{\mathrm{int}}=(H-1)N^2,\qquad
 \Omega=\omega+D_{\mathrm{int}}.
\]
The independent Gaussian base now has variances
\begin{equation}
 \operatorname{Var}_0(W^1_{ji})=\frac{N^{1-\gamma}}d,\qquad
 \operatorname{Var}_0(W^\ell_{ji})=N^{-\gamma}\ (\ell\ge2),\qquad
 \operatorname{Var}_0(v_j)=N^{-\gamma-1}.
 \label{eq:usual-base-scales}
\end{equation}
For this section redefine
\[
 u_N=\frac dN\|W^1\|_F^2,\qquad q_N=N\|v\|_2^2,\qquad Q_N=G(u_N+q_N).
\]
The prior is
\begin{equation}
 \pi^{\mathrm{us}}_{N,\gamma}(dw)\ \propto\
 \exp\left[-\frac\omega2(G-1)(u_N+q_N)
            -\frac\Omega2(G^2-1)\right]\pi^0_{N,\gamma}(dw).
 \label{eq:usual-soft-prior}
\end{equation}
Thus only a soft gain regularizer is added to~\eqref{eq:usual-base-scales}. Its coefficient $\Omega$ pays for the internal parameter volume. This coefficient depends on the learner architecture, not on the teacher.

The conditional endpoint Gaussian variances are the base variances divided by $G$. In particular
\begin{equation}
 \omega Q_N\sim\chi^2_{N(d+1)},\qquad
 \frac{N(d+1)}\omega\le d+1,\qquad
 P(w)\le\frac{Q_N}{2\sqrt d}\le\frac{Q_N}{2}.
 \label{eq:usual-energy}
\end{equation}
The same radial normalizer argument as Lemma~\ref{lem:gain-moment}, with $\Omega$ in place of $\omega$, gives
\begin{equation}
 \E_\pi G^2\le1+\frac{D_{\mathrm{int}}}{\Omega}\le2.
 \label{eq:usual-gain}
\end{equation}
Conditional Gaussian cancellation also gives $\E_\pi z_w(x)^2\le N/(d\omega^2)\le1$, so the free-energy bound remains $F_S\le2$.

\begin{lemma}[Selective block transport]
\label{lem:selective-block}
Fix a teacher of width $C$, depth $H$, and path variation $P$. Put
\[
 \overline P=\sqrt d\,P,\qquad
 D_C=d+1+2(H-1)C.
\]
For every $N\ge C$ and $0<\eta<1$, there is a comparison law $\rho_{g,\eta}$ for~\eqref{eq:usual-soft-prior} satisfying
\begin{align}
 \frac2\omega\KL(\rho_{g,\eta}\|\pi^{\mathrm{us}}_{N,\gamma})
 &\le3\overline P+(H-1)C+2D_C\log(1/\eta),
 \label{eq:usual-entropy}\\
 \left(\int\|z_w-g\|_{\infty,\mathcal X}^2d\rho_{g,\eta}\right)^{1/2}
 &\le6(H+1)(\overline P+1)\eta.
 \label{eq:usual-tube}
\end{align}
\end{lemma}
\begin{proof}
Apply flow balancing and zero-pad the resulting width-$C$ teacher into the first $C$ coordinates of the width-$N$ network. Its internal active blocks $A^\ell$ have absolute operator norm at most one and squared Frobenius norm at most $C$. Multiply the first matrix by $\sqrt N/d^{1/4}$ and the readout by its inverse. Positive homogeneity preserves the function and makes the raw endpoint energies
$u_N=q_N=\overline P$.

Draw the whole internal tuple $X$ from its actual prior marginal. For each layer, in the decomposition into active and inactive coordinates define
\[
 B^\ell(X)=
 \begin{pmatrix} A^\ell&0\\0&X^\ell_{\mathrm{inactive,inactive}}\end{pmatrix},
 \qquad T_\eta X^\ell=(1-\eta)B^\ell(X)+\eta X^\ell.
\]
The inactive block is unchanged. Exactly $2CN-C^2$ coordinates in each internal layer are multiplied by $\eta$, with a translation on its active block. The full affine Jacobian is therefore
$\eta^{(H-1)(2CN-C^2)}$.

A submatrix of a nonnegative matrix has operator norm no greater than the full matrix. Consequently
\[
 \bigl\||B^\ell(X)|\bigr\|_{\op}
 =\max\{\||A^\ell|\|_{\op},
          \||X^\ell_{\mathrm{inactive,inactive}}|\|_{\op}\}
 \le g_\ell(X).
\]
Convexity gives $g_\ell(T_\eta X)\le g_\ell(X)$ and hence $G(T_\eta X)\le G(X)$. The internal density has the same form as~\eqref{eq:internal-marginal}, with quadratic precision $N^\gamma$ and gain-penalty coefficient $\Omega$. In the exact log-density ratio, both gain-potential differences are nonpositive. The unchanged Gaussian coordinates cancel, and the changed coordinates have centered prior means. Their normalized quadratic mean cost is at most $(H-1)C$.

Conditional on the transformed internals, use endpoint Gaussians with the teacher means and $\eta^2$ times the actual conditional covariances. Their normalized mean cost is at most $2\overline P\E G(T_\eta X)\le3\overline P$. All covariance entropy terms together cost at most
\[
 \frac{2\{N(d+1)+(H-1)(2CN-C^2)\}}{N^\gamma}\log(1/\eta)
 \le2D_C\log(1/\eta).
\]
This proves~\eqref{eq:usual-entropy}.

The random block-diagonal baseline $B(X)$ computes exactly the teacher: the first layer and readout have zero inactive coordinates, and ReLU maps zero to zero. Its teacher hidden states stay in the active block and have Euclidean norm divided by $\sqrt N$ at most $\sqrt{\overline P}$. Moreover
\[
 \|T_\eta X^\ell-B^\ell(X)\|_{\op}
 \le2\eta g_\ell(X).
\]
Telescoping with this baseline gives exactly the same internal perturbation bound as Lemma~\ref{lem:comparison}, using~\eqref{eq:usual-gain}; independence of internal layers is unnecessary. The conditional first-layer noise has normalized squared Frobenius mean $\eta^2N/(\omega G_T)\le\eta^2/G_T$, where $G_T=G(T_\eta X)$. The conditional readout-noise bound is at most
$\eta^2(N/\omega)\overline P\E G_T+\eta^4N^2/\omega^2\le\eta^2(2\overline P+1)$.
The same Minkowski estimate therefore proves~\eqref{eq:usual-tube}.
\end{proof}

\begin{theorem}[A path oracle with the usual Gaussian base]
\label{thm:usual-soft-oracle}
For the prior~\eqref{eq:usual-soft-prior}, use the Gaussian likelihood with observation variance $n\lambda/\omega$ and the free-energy root~\eqref{eq:calibration}. In~\eqref{eq:stat-constants} take $D_0=d+1$. For a comparison teacher define
\[
 \eta_n=\frac1{64n(H+1)(\overline P+1)^2},\qquad
 T_g=6+3\overline P+(H-1)C+2D_C\log(1/\eta_n),\qquad
 \mathcal B_g=2T_g+K_n+1+c_0.
\]
Then the exact population oracle~\eqref{eq:population-oracle}, with these constants, holds simultaneously at every $N\ge C$ and every $\gamma\ge1$. In particular a sufficient sample count for excess at most $\varepsilon$ is
\begin{equation}
 n=\widetilde O\left(
 H\bigl[1+\sqrt d\,P+HC+d\bigr]^2
 \left\{\frac{R_*}{\varepsilon^2}+\frac1\varepsilon\right\}\right).
 \label{eq:usual-soft-samples}
\end{equation}
The prior and the selected likelihood use none of $C$, $P$, or $R_*$. At fixed $d$ and bounded teacher gains and endpoint norms, $P=O(C)$, so the bound is $\widetilde O(C^2H^3(R_*/\varepsilon^2+1/\varepsilon))$.
\end{theorem}
\begin{proof}
Equations~\eqref{eq:usual-energy} and~\eqref{eq:usual-gain} provide the amplitude, capacity, and endpoint-moment assumptions used in the certificate proof, with $D_0=d+1$. The prior risk is at most two. Lemma~\ref{lem:selective-block} at $\eta_n$ gives
$F_S(\lambda)\le\widehat R_S(g)+1/(2n)+\lambda T_g$.
For the population conversion use $E=6+3\overline P+(H-1)C\le T_g$, so $|g|\le P\le E$. Every remaining step and the numerical constants 25 and 280 in Theorem~\ref{thm:main} are unchanged. The displayed sample rate follows from $a_n^2=\widetilde O(H/n)$ and $\mathcal B_g=\widetilde O(1+\sqrt dP+HC+d)$.
\end{proof}

\section{Ordinary spectral gains: a more permissive prior in fixed input dimension}
\label{sec:ordinary-spectral-variant}
Absolute matrix norms control path variation but discard cancellation between positive and negative entries. A closely related prior can use ordinary spectral norms instead. Its guarantee depends more strongly on input dimension, while its dependence on depth is only logarithmic once the teacher's spectral product is controlled. Ordinary spectral gains are the quantities constrained by spectral normalization in neural-network training~\cite{spectral-normalization}; the theorem below is a Bayesian guarantee for the explicitly stated prior, not a claim that ordinary training automatically yields a bounded gain product.

Define
\[
 G_{\mathrm{sp}}(W)=\prod_{\ell=2}^H\max\{1,\|W^\ell\|_{\op}\}.
\]
Replace $G$ by $G_{\mathrm{sp}}$ in the prior~\eqref{eq:prior}, leaving the Gaussian base scales and $\omega=N(d+1)+(H-1)N^{\gamma+1}$ unchanged. Denote this full-support prior by $\pi_{\mathrm{sp}}$ and its endpoint energy by
\[
 Q_{\mathrm{sp}}=G_{\mathrm{sp}}(u_N+q_N).
\]
For a ReLU implementation define its spectral product complexity
\begin{equation}
 S(w)=\|W^1\|_F\,\|v\|_2\prod_{\ell=2}^H\|W^\ell\|_{\op}.
 \label{eq:spectral-product}
\end{equation}
Every represented function vanishes at the origin and has Lipschitz constant and unit-ball amplitude at most $S(w)\le Q_{\mathrm{sp}}/2$. The origin need not belong to the actual input domain; the network extends to the whole ambient unit ball.

\begin{lemma}[Spectral-product comparison]
\label{lem:spectral-comparison}
For every width-$C$, depth-$H$ ReLU teacher $g$ with spectral product $S$, every $N\ge C$, every $\gamma\ge1$, and $0<\eta<1$, a comparison law for $\pi_{\mathrm{sp}}$ satisfies
\[
 \frac2\omega\KL(\rho\|\pi_{\mathrm{sp}})
 \le3S+1+2\log(1/\eta),\qquad
 \left(\E_\rho\|z_w-g\|_{\infty,\mathcal X}^2\right)^{1/2}
 \le6(H+1)(S+1)\eta.
\]
Moreover $\omega Q_{\mathrm{sp}}$ has the chi-square law with $N(d+1)$ degrees of freedom under the prior, and its internal gain obeys $\E G_{\mathrm{sp}}^2\le2$.
\end{lemma}
\begin{proof}
If $g$ is zero, use its all-zero implementation. Otherwise all factors in~\eqref{eq:spectral-product} are positive. Divide each internal matrix by its spectral norm and multiply the readout by the product of those norms. ReLU homogeneity preserves the function and leaves every internal norm equal to one. Clone neurons into $k=\lfloor N/C\rfloor$ equal groups, leaving the remaining coordinates zero. The internal clone of a matrix $A$ is $A\otimes(J_k/k)$, so its spectral norm is unchanged and its squared Frobenius norm is at most $C$. The cloned endpoint energies have product $S^2$. Scaling the first layer by a positive scalar and the readout by its reciprocal therefore makes both endpoint energies equal to $S$, again preserving the function.

The proof of Lemma~\ref{lem:gain-moment} uses only positive homogeneity and monotonicity of the factors $\max(1,\|W\|)$ under scalar contraction. It applies to $G_{\mathrm{sp}}$ without a change. Conditional endpoint covariance normalization gives the same chi-square identity as before.

For the comparison, transport the actual internal tuple toward the normalized cloned teacher by $T_\eta X=(1-\eta)A+\eta X$. Convexity of ordinary spectral norm gives
$G_{\mathrm{sp}}(T_\eta X)\le1+\eta(G_{\mathrm{sp}}(X)-1)\le G_{\mathrm{sp}}(X)$.
Thus both internal gain-potential differences in the exact log-density ratio are nonpositive. The teacher internal Frobenius cost is at most $C/N\le1$ after normalization, and the endpoint mean cost is at most $2\sqrt2S\le3S$. All dimension factors are unchanged, proving the entropy bound. Finally the uniform-error proof of Lemma~\ref{lem:comparison} uses only ordinary operator norms, endpoint energies, and the moment $\E G^2\le2$. The same telescoping calculation therefore gives the displayed tube estimate with $S$ in place of $P$.
\end{proof}

\begin{lemma}[Capacity of a Lipschitz ball]
\label{lem:Lipschitz-capacity}
Let $s_d=\max\{d,2\}$. There is a finite constant $C_d$, depending only on $d$, such that the class $\mathcal L_t$ of real functions on the Euclidean unit ball with $f(0)=0$ and Lipschitz constant at most $t$ satisfies
\[
 \Rad_n(\mathcal L_t)\le C_dt\,\log^2(en)\,n^{-1/s_d}
\]
on every sample. In particular the same bound holds for the union over all widths of $Q_{\mathrm{sp}}\le t$ networks, with an inessential factor of two of slack.
\end{lemma}
\begin{proof}
Scale to $t=1$. Such functions are bounded by one. An $\varepsilon/4$-net of the input ball has at most $(C\sqrt d/\varepsilon)^d$ points. Quantizing function values on that net to multiples of $\varepsilon/2$ gives a sup-norm cover with
\[
 \log\mathcal N(\varepsilon,\mathcal L_1,\|\cdot\|_\infty)
 \le(C\sqrt d/\varepsilon)^d\log(C/\varepsilon)
 \qquad(0<\varepsilon<1).
\]
Indeed two functions with the same quantized values differ by at most $\varepsilon$ everywhere; choose one function from each nonempty code class. The truncated Rademacher entropy integral then bounds the complexity by
\[
 C_d\inf_{0<r<1}\left\{r+\frac1{\sqrt n}
        \int_r^1\varepsilon^{-d/2}\sqrt{\log(C/\varepsilon)}\,d\varepsilon\right\}.
\]
For $d=1$ the integral converges as $r$ tends to zero. For $d=2$, take $r=n^{-1/2}$, giving at most $C_2n^{-1/2}\log^{3/2}(en)$. For $d\ge3$, take $r=n^{-1/d}$, giving at most $C_dn^{-1/d}\sqrt{\log(en)}$. The displayed common bound follows, including small $n$ after increasing $C_d$. The entropy-integral inequality itself follows by chaining successive covers at dyadic scales, applying the finite-class Rademacher bound to their increments, and stopping at radius $r$; the remaining increments have empirical norm at most $r$.
\end{proof}

\begin{theorem}[A spectral-product oracle in fixed input dimension]
\label{thm:spectral-oracle}
Assume any IID law on $\mathcal X\times[-1,1]$. For the prior $\pi_{\mathrm{sp}}$, use the same observed free-energy root~\eqref{eq:calibration}, with $D_0=1$ and with
\[
 a_{n,d}=\max\left\{n^{-1/2},\
 2C_*\left[\sqrt2 C_d\log^2(en)n^{-1/s_d}\log^{3/2}(2n)
                 +2\sqrt{t_n/n}\right]\right\}
\]
in place of $a_n$. Here $t_n$ and all other constants remain as in~\eqref{eq:stat-constants}. The ordinary raw posterior mean then satisfies, for every fixed width-$C$ teacher $g$ with spectral product $S$, with probability at least $1-\delta$, simultaneously for all $N\ge C$ and $\gamma\ge1$,
\[
 R(\widehat f)\le R(g)+25a_{n,d}\mathcal B_g\sqrt{R(g)}
                         +280a_{n,d}^2\mathcal B_g^2+n^{-4},
\]
where
\[
 \eta_n=[64n(H+1)(S+1)^2]^{-1},\qquad
 T_g=6+3S+2\log(1/\eta_n),\qquad
 \mathcal B_g=2T_g+K_n+1+c_0.
\]
The learner uses neither the teacher resources nor its risk.
\end{theorem}
\begin{proof}
The pointwise amplitude bound is $|z|\le Q_{\mathrm{sp}}/2$, and Lemma~\ref{lem:Lipschitz-capacity} replaces Lemma~\ref{lem:capacity}. Apply the smooth-loss inequality and dyadic energy peeling exactly as in Lemma~\ref{lem:stat-event}; the stated $a_{n,d}$ gives its identical conclusion. The chi-square energy law, normalized endpoint dimension, and prior second-moment bound are unchanged, so the weighted empirical-objective tail and entropy lemmas apply verbatim. The root argument therefore gives the same observable certificate and empirical oracle. Lemma~\ref{lem:spectral-comparison} supplies the free-energy comparison with $S$ in place of $P$. Finally $|g|\le S$ gives the same fixed-comparator Bernstein step and the stated constants 25 and 280.
\end{proof}

At fixed input dimension, a sufficient sample count is
\begin{equation}
 n=\widetilde O_d\left((1+S)^{s_d}
       \left\{\frac{R_*^{s_d/2}}{\varepsilon^{s_d}}
                        +\frac1{\varepsilon^{s_d/2}}\right\}\right).
 \label{eq:spectral-samples}
\end{equation}
Depth appears only through logarithms once $S$ is held fixed. If first-row Euclidean norms are at most $B_1$, readout coefficients are at most $B_2$, and the product of ordinary internal spectral norms is at most one, then $S\le CB_1B_2$. In dimensions one and two the result is therefore quadratic in teacher width, up to logarithms, without polynomial depth cost. In larger fixed dimension it remains polynomial, with width exponent $d$ and the displayed accuracy powers.

This variant does not give a polynomial jointly in input dimension and teacher width: both the exponent and the constants in~\eqref{eq:spectral-samples} deteriorate with dimension. Its purpose is different. It replaces absolute spectral gain by the ordinary Lipschitz amplification of the hidden layers, retains cancellation between signed weights, and shows precisely what the more permissive prior can guarantee without the path-variation capacity argument.

\section{Scope and interpretation}
\paragraph{A full-support prior with an explicit inductive preference.}
Every finite parameter vector has positive prior density. Large internal gains reduce endpoint variances instead of being forbidden. Theorem~\ref{thm:prior} describes the resulting probabilities of function neighborhoods. Theorem~\ref{thm:energy-characterization} identifies the limiting minimum normalized energy with twice minimal path variation. The conditional Gaussian normalizers, including the tilted internal marginal, are included in the probability proof.

\paragraph{Why path variation is a natural quantity.}
Positive neuron rescalings leave a ReLU function and its path products unchanged. Path regularizers and this rescaling symmetry are established tools in neural-network optimization~\cite{path-sgd}. Total path variation also controls width-independent statistical complexity~\cite{barron-klusowski}. The mixed version here uses Euclidean first-row norms because the input domain is a Euclidean unit ball. Our contribution is the explicit Gaussian construction whose function probabilities and ordinary posterior satisfy the stated bounds; we do not claim that the path complexity measure itself is new.

\paragraph{What changed from independent Gaussian weights.}
The base internal coordinates have variance $N^{-\gamma}$, and the base endpoint variances are $N/\omega$ and $1/(\omega N)$. The soft gain penalty changes the internal marginal and gives gain-normalized conditional endpoint variances. At $\gamma=1$, $\omega$ is exactly the Gaussian parameter count. No additional width threshold or logarithmic precision correction is needed. The theorem does not prove the same path-variation bound for the unchanged independent-coordinate prior.

\paragraph{Other internal Gaussian scales.}
The proof also allows arbitrary positive internal precisions $\alpha_\ell$. Use base internal variance $1/\alpha_\ell$ and set
$\omega=N(d+1)+N\sum_{\ell=2}^H\max\{N,\alpha_\ell\}$.
The gain moment still has $D_{\mathrm{int}}/\omega\le1$, the total dimension ratio is at most one, and the cloned internal mean cost is at most $C/N$. Every other argument is unchanged. This is a calibrated family of softly regularized priors, not an assertion about arbitrary unchanged independent Gaussian exponent tuples.

\paragraph{Small width versus small path variation.}
Every finite ReLU function has a flow-balanced representation at its existing width and depth. Nevertheless a small network can have large path variation, and the theorem does not give unrestricted polynomial dependence on width and depth. The quadratic-width, linear-depth corollary requires bounded endpoint norms and bounded product of internal absolute spectral gains. No such teacher bound is used to choose the prior or likelihood.

\paragraph{Noise, computation, and verification.}
The labels may follow any bounded IID law, without a network target or Gaussian noise assumption. The data-dependent likelihood yields faster rates at smaller comparison risk and predicts by averaging raw network outputs. The theorem provides neither efficient evaluation of free energy nor a gradient-descent guarantee. These conventional proofs have not been formalized in Lean.

\begin{thebibliography}{9}
\bibitem{sst} Nathan Srebro, Karthik Sridharan and Ambuj Tewari. \emph{Smoothness, Low Noise and Fast Rates}. Advances in Neural Information Processing Systems 23, 2010. \url{https://www.cs.cornell.edu/~sridharan/smoothness.pdf}.
\bibitem{path-sgd} Behnam Neyshabur, Ruslan Salakhutdinov and Nathan Srebro. \emph{Path-SGD: Path-Normalized Optimization in Deep Neural Networks}. Advances in Neural Information Processing Systems 28, 2015. \url{https://arxiv.org/abs/1506.02617}.
\bibitem{barron-klusowski} Andrew R. Barron and Jason M. Klusowski. \emph{Complexity, Statistical Risk, and Metric Entropy of Deep Nets Using Total Path Variation}. 2019. \url{https://arxiv.org/abs/1902.00800}.
\bibitem{spectral-normalization} Takeru Miyato, Toshiki Kataoka, Masanori Koyama and Yuichi Yoshida. \emph{Spectral Normalization for Generative Adversarial Networks}. International Conference on Learning Representations, 2018. \url{https://arxiv.org/abs/1802.05957}.
\end{thebibliography}
\end{document}
